\section{Complete elliptic integrals}\label{app:elliptic}

This appendix derives every property of $K$ and $E$ used in the paper. The
modulus is $k\in[0,1)$ and the complementary modulus is $k'=\sqrt{1-k^2}$;
we write $K'=K(k')$ and $E'=E(k')$.

\subsection{Power series}

For $\abs u<1$, $(1-u)^{-1/2}=\sum_{n\ge0}c_nu^n$ and
$(1-u)^{1/2}=\sum_{n\ge0}c_nu^n/(1-2n)$, with $c_n=\binom{2n}n4^{-n}$; the
first is the binomial series, $\binom{-1/2}n(-1)^n=c_n$, and the second follows
by integrating the first term by term. The Wallis integral
$\frac2\pi\int_0^{\pi/2}\sin^{2n}\phi\dd\phi=c_n$ follows by induction from
the reduction formula
$\int_0^{\pi/2}\sin^{2n}=\frac{2n-1}{2n}\int_0^{\pi/2}\sin^{2n-2}$ and the
recurrence $c_n=\frac{2n-1}{2n}c_{n-1}$. With $u=k^2\sin^2\phi$ and uniform
convergence for $k<1$,
\begin{equation}\label{eq:KEseries}
K(k)=\frac\pi2\sum_{n\ge0}c_n^2k^{2n},\qquad
E(k)=\frac\pi2\sum_{n\ge0}\frac{c_n^2}{1-2n}k^{2n}.
\end{equation}
Subtracting, $K-E=\frac\pi2\sum_{n\ge1}c_n^2\frac{2n}{2n-1}k^{2n}$, which is
the series of \cref{prop:Lambda}(i).\footnote{The coefficients
$d_n=c_n^2\,2n/(2n-1)$ are $\frac12,\frac3{16},\frac{15}{128},\frac{175}{2048},
\dots$; the ratio $d_{n+1}/d_n=(2n+1)(2n-1)/((2n+2)2n)$ follows from
$c_{n+1}/c_n=(2n+1)/(2n+2)$.}

\subsection{Derivatives}

Differentiating under the integral sign,
\begin{equation}\label{eq:derivs}
\frac{\dd E}{\dd k}=\frac{E-K}k,\qquad
\frac{\dd K}{\dd k}=\frac{E-k'^2K}{k\,k'^2}.
\end{equation}
The first is immediate: $\frac{\dd}{\dd k}\sqrt{1-k^2\sin^2\phi}
=-k\sin^2\phi/\sqrt{\cdot}$, and $-k^2\sin^2\phi=(1-k^2\sin^2\phi)-1$. For the
second, the derivative is $\int_0^{\pi/2}k\sin^2\phi\,(1-k^2\sin^2\phi)^{-3/2}
\dd\phi$, and the identity
\begin{multline*}
\frac{\dd}{\dd\phi}\frac{\sin\phi\cos\phi}{\sqrt{1-k^2\sin^2\phi}}\\
=\frac1{k^2}\Big(\sqrt{1-k^2\sin^2\phi}-\frac{k'^2}{(1-k^2\sin^2\phi)^{3/2}}\Big),
\end{multline*}
integrated over $[0,\pi/2]$, gives $\int(1-k^2\sin^2)^{-3/2}=E/k'^2$; then
$k^2\sin^2\phi(1-k^2\sin^2\phi)^{-3/2}=(1-k^2\sin^2)^{-3/2}-(1-k^2\sin^2)^{-1/2}$
yields \eqref{eq:derivs}.\footnote{To check the displayed identity, put
$s=\sin\phi$, $c=\cos\phi$, $\Delta=1-k^2s^2$. The derivative of $sc\,\Delta^{-1/2}$
is $(c^2-s^2)\Delta^{-1/2}+k^2s^2c^2\Delta^{-3/2}=\Delta^{-3/2}\big((c^2-s^2)\Delta
+k^2s^2c^2\big)$, and with $c^2=1-s^2$ the bracket is
$1-2s^2+k^2s^4=\frac1{k^2}\big(\Delta^2-k'^2\big)$.}

\subsection{Legendre's relation}

\begin{lemma}\label{lem:legendre}
$E(k)K'+E'K(k)-K(k)K'=\pi/2$ for $0<k<1$.
\end{lemma}

\begin{proof}
Let $L(k)$ be the left side. Since $\dd k'/\dd k=-k/k'$, \eqref{eq:derivs}
applied at $k'$ gives
\[
\frac{\dd E'}{\dd k}=\frac{k(K'-E')}{k'^2},\qquad
\frac{\dd K'}{\dd k}=-\frac{E'-k^2K'}{kk'^2}.
\]
Multiply $\dd L/\dd k$ by $kk'^2$ and insert the four derivatives:
\begin{multline*}
kk'^2L'=k'^2(E-K)K'-E(E'-k^2K')+k^2K(K'-E')\\
+E'(E-k'^2K)-K'(E-k'^2K)+K(E'-k^2K').
\end{multline*}
Expanding, the terms in $KK'$ and in $EE'$ cancel in pairs, and what
remains is
\[
kk'^2L'=(k^2+k'^2-1)\,(EK'-KE')=0 .
\]
So $L$ is constant. As $k\to0$, $K,E\to\pi/2$ and $E'\to1$, while
$K-E=O(k^2)$ by \eqref{eq:KEseries} and $K'=O(\log(1/k))$ by
\eqref{eq:Klog} below; hence $L=E'K+(E-K)K'\to\pi/2$.
\end{proof}

\subsection{The arithmetic--geometric mean and Landen's transformation}

For $p,q>0$ let
\[
I(p,q)=\int_0^{\pi/2}\frac{\dd\phi}{\sqrt{p^2\cos^2\phi+q^2\sin^2\phi}},
\]
so that $K(k)=I(1,k')$.

\begin{lemma}[Gauss]\label{lem:gauss}
$I(p,q)=I\big(\frac{p+q}2,\sqrt{pq}\big)$. Consequently
$I(p,q)=\pi/(2M(p,q))$, with $M$ the arithmetic--geometric mean of $p$
and $q$.
\end{lemma}

\begin{proof}
With $x=q\tan\phi$, $I(p,q)=\frac12\int_\R\big((p^2+x^2)(q^2+x^2)\big)^{-1/2}\dd x$.
Substitute $x=\frac12(t-pq/t)$, $t>0$, which maps $(0,\infty)$ increasingly
onto $\R$, with $\dd x=\frac12(1+pq/t^2)\dd t$. Then
$4x^2+(p+q)^2=(t^2+p^2)(t^2+q^2)/t^2$ and $x^2+pq=\frac14(t+pq/t)^2$, so
\[
\big((p_1^2+x^2)(q_1^2+x^2)\big)^{-1/2}\dd x
=\frac{2\,\dd t}{\sqrt{(t^2+p^2)(t^2+q^2)}},
\]
with $p_1=\frac{p+q}2$, $q_1=\sqrt{pq}$. Integrating over $t\in(0,\infty)$
and using that the integrand is even in $t$,
\begin{align*}
2I(p_1,q_1)&=2\int_0^\infty\frac{\dd t}{\sqrt{(t^2+p^2)(t^2+q^2)}}\\
&=\int_\R\frac{\dd t}{\sqrt{(t^2+p^2)(t^2+q^2)}}=2I(p,q).
\end{align*}
The
iterates converge to $M(p,q)$, and $I(M,M)=\pi/(2M)$.
\end{proof}

\begin{proposition}[Landen]\label{prop:landen}
Let $0<k'<1$ and $\kappa=(1-k')/(1+k')$. Then
\[
K(k)=(1+\kappa)K(\kappa),\qquad E(k)=(1+k')E(\kappa)-k'K(k).
\]
\end{proposition}

\begin{proof}
By \cref{lem:gauss}, $K(k)=I(1,k')=I\big(\frac{1+k'}2,\sqrt{k'}\big)
=\frac2{1+k'}I\big(1,\frac{2\sqrt{k'}}{1+k'}\big)$, and
$\frac{2\sqrt{k'}}{1+k'}=\sqrt{1-\kappa^2}=\kappa'$, so
$K(k)=\frac2{1+k'}K(\kappa)=(1+\kappa)K(\kappa)$. For $E$, regard $k$ as a
function of $\kappa$, $k=2\sqrt\kappa/(1+\kappa)$, $k'=(1-\kappa)/(1+\kappa)$,
and differentiate the relation for $K$ with \eqref{eq:derivs}. Since
\begin{gather*}
\frac{\dd k}{\dd\kappa}=\frac{1-\kappa}{\sqrt\kappa(1+\kappa)^2},\qquad
kk'^2=\frac{2\sqrt\kappa(1-\kappa)^2}{(1+\kappa)^3},\\
\text{so}\qquad
\frac{\dd k/\dd\kappa}{kk'^2}=\frac{1+\kappa}{2\kappa(1-\kappa)},
\end{gather*}
the derivative of $K(k)$ is $\big(E(k)-k'^2K(k)\big)\frac{1+\kappa}{2\kappa(1-\kappa)}$,
and that of $(1+\kappa)K(\kappa)$ is
$K(\kappa)+\big(E(\kappa)-(1-\kappa^2)K(\kappa)\big)/(\kappa(1-\kappa))$.
Equating, multiplying by $2\kappa(1-\kappa)/(1+\kappa)$ and using
$k'^2(1+\kappa)=(1-\kappa)^2/(1+\kappa)$,
\begin{align*}
E(k)&=\frac{2E(\kappa)}{1+\kappa}+\frac{K(\kappa)}{1+\kappa}\Big((1-\kappa)^2
+2\kappa(1-\kappa)\\
&\qquad-2(1-\kappa)(1+\kappa)\Big)\\
&=\frac{2E(\kappa)}{1+\kappa}-(1-\kappa)K(\kappa).
\end{align*}
Since $2/(1+\kappa)=1+k'$ and $(1-\kappa)K(\kappa)=k'K(k)$, this is the
claim.
\end{proof}

\subsection{Behaviour at the ends}

As $k\to0$, \eqref{eq:KEseries} gives $K,E=\frac\pi2+O(k^2)$. As $k\to1$,
$E(k)\to1$ and
\begin{equation}\label{eq:Klog}
K(k)=\log\frac4{k'}+O\big(k'^2\log k'\big).
\end{equation}
The limit of $E$ is dominated convergence. For $K$, the substitution
$x=k'\tan\phi$ of \cref{lem:gauss} gives
$K(k)=\int_0^\infty\big((1+x^2)(k'^2+x^2)\big)^{-1/2}\dd x$. Split at $x=1$:
\[
\int_0^1\frac{\dd x}{\sqrt{k'^2+x^2}}=\operatorname{arsinh}\frac1{k'}
=\log\frac2{k'}+o(1),
\]
and, by dominated convergence,
\begin{gather*}
\int_0^1\frac{(1+x^2)^{-1/2}-1}{\sqrt{k'^2+x^2}}\dd x\to\log\frac2{1+\sqrt2},\\
\int_1^\infty\frac{\dd x}{\sqrt{(1+x^2)(k'^2+x^2)}}\to\log(1+\sqrt2),
\end{gather*}
the first limit being
$\int_0^1\big(\frac1{x\sqrt{1+x^2}}-\frac1x\big)\dd x=\log\frac2{1+\sqrt2}$,
both from the primitive $-\operatorname{arsinh}(1/x)$ of
$1/(x\sqrt{1+x^2})$. The three limits add up to $\log(4/k')$.
The error term is \cite[19.12.1]{dlmf}.

\section{The elementary integrals}\label{app:integrals}

\subsection{Moments of the arcsine weight}

For $A<B$ put $m=\frac{A+B}2$, $d=\frac{B-A}2$ and
$\sigma=m-d\cos\phi$, $0\le\phi\le\pi$. Then
$(\sigma-A)(B-\sigma)=d^2\sin^2\phi$ and $\dd\sigma=d\sin\phi\dd\phi$, so
$\dd\sigma/\sqrt{(\sigma-A)(B-\sigma)}=\dd\phi$, and
\begin{align*}
\int_A^B\frac{\dd\sigma}{\sqrt{(\sigma-A)(B-\sigma)}}&=\pi,\\
\int_A^B\frac{\sigma\dd\sigma}{\sqrt\cdot}&=\pi m=\frac\pi2(A+B),\\
\int_A^B\frac{\sigma^2\dd\sigma}{\sqrt\cdot}&=\pi\Big(m^2+\frac{d^2}2\Big)\\
&=\frac\pi8\big(3A^2+2AB+3B^2\big),\\
\int_A^B\frac{\dd\sigma}{\sigma\sqrt\cdot}&=\int_0^\pi\frac{\dd\phi}{m-d\cos\phi}\\
&=\frac\pi{\sqrt{m^2-d^2}}=\frac\pi{\sqrt{AB}} .
\end{align*}
The last integral is the classical $\int_0^\pi\dd\phi/(m-d\cos\phi)=\pi/\sqrt{m^2-d^2}$
for $m>\abs d$, which follows from $u=\tan(\phi/2)$.\footnote{With
$u=\tan(\phi/2)$, $\cos\phi=(1-u^2)/(1+u^2)$ and $\dd\phi=2\dd u/(1+u^2)$, so
the integral is $\int_0^\infty2\dd u/\big((m-d)+(m+d)u^2\big)
=\pi/\sqrt{(m-d)(m+d)}$.}

These give \cref{lem:moments}(iii),(iv): with
$x\,p(x)\dd x=\frac12\varpi(\sigma)\dd\sigma$ and
$p(x)\dd x/x=\frac12\varpi(\sigma)\dd\sigma/\sigma$,
\begin{align*}
\int_a^b\frac{p}x\dd x&=\frac12\Big((A+B)\frac\pi{\sqrt{AB}}-2\pi\Big)\\
&=\frac{\pi(b-a)^2}{2ab},\\
\int_a^bx^3p\dd x&=\frac12\Big(\frac\pi2(A+B)^2\\
&\qquad-\frac\pi4(3A^2+2AB+3B^2)\Big)\\
&=-\frac\pi8(B-A)^2 .
\end{align*}

\subsection{The integrals $I_0$, $I_2$ and $J$}

With $t^2=b^2-(b^2-a^2)\sin^2\phi$ (\cref{prop:J}),
$t\dd t=-(b^2-a^2)\sin\phi\cos\phi\dd\phi$ and $r(t)=(b^2-a^2)\sin\phi\cos\phi$,
so $\dd t/r=-\dd\phi/t$ and
\begin{gather*}
I_0=\int_0^{\pi/2}\frac{\dd\phi}{b\sqrt{1-k^2\sin^2\phi}}=\frac{K(k)}b,\\
I_2=\int_0^{\pi/2}b\sqrt{1-k^2\sin^2\phi}\dd\phi=bE(k).
\end{gather*}
Hence $J=(a^2+b^2)I_0-2I_2=b\big((1+k'^2)K(k)-2E(k)\big)$. By
\cref{prop:landen}, $-2E(k)=-2(1+k')E(\kappa)+2k'K(k)$, and
$(1+k'^2)+2k'=(1+k')^2$, so
\begin{align*}
J&=b(1+k')^2K(k)-2b(1+k')E(\kappa)\\
&=b(1+k')^2(1+\kappa)K(\kappa)-2b(1+k')E(\kappa).
\end{align*}
With $1+k'=(a+b)/b$ and $1+\kappa=2b/(a+b)$ the coefficients are $2(a+b)$
and $2(a+b)$, which is \eqref{eq:J}.

\subsection{The integrals on $(0,a)$}

With $s=a\sin\phi$,
\[
\sqrt{(a^2-s^2)(b^2-s^2)}=a\cos\phi\,b\sqrt{1-k'^2\sin^2\phi},
\]
so
\begin{gather*}
\int_0^a\frac{\dd s}{\sqrt{(a^2-s^2)(b^2-s^2)}}=\frac{K(k')}b,\\
\int_0^a\frac{s^2\dd s}{\sqrt{(a^2-s^2)(b^2-s^2)}}=b\big(K(k')-E(k')\big),
\end{gather*}
the second because $a^2\sin^2\phi=b^2k'^2\sin^2\phi$ and
\[
\int_0^{\pi/2}\frac{k'^2\sin^2\phi\,\dd\phi}{\sqrt{1-k'^2\sin^2\phi}}=K(k')-E(k').
\]
These are
the integrals in the proof of \cref{thm:extremal}(ii), where
\[
\frac{m^2}bK'-b(K'-E')=\frac b{K}\big(EK'-KK'+E'K\big)=\frac{\pi b}{2K}
\]
by $m^2=b^2E/K$ and \cref{lem:legendre}.

\section{The extremal profile in closed form}\label{app:profiles}

\subsection{Three formulas for $\alpha_\star$}

\begin{lemma}\label{lem:alphaforms}
For $a<x<b$,
\begin{align*}
\alpha_\star(x)&=\tau_\star\int_a^x\frac{m^2-s^2}{r(s)}\dd s\\
&=\frac{2x\,r(x)\theta}\pi\PV\int_a^b\frac{\dd t}{r(t)(x^2-t^2)}\\
&=\frac{2\theta\,r(x)}\pi\int_a^b\frac{\dd t}{r(t)(x+t)} .
\end{align*}
\end{lemma}

The last form shows at once that $\alpha_\star>0$ on $(a,b)$ and that it
vanishes like $\sqrt{x-a}$ and $\sqrt{b-x}$ at the band edges.

\begin{proof}
The first equality is \cref{thm:extremal}(ii), the second
\cref{prop:profile}(ii),(iv) with $g=\theta$. For the third, apply Cauchy's
formula to $1/R$, which is analytic off $S$, $O(z^{-2})$ at infinity, and
has the jumps $1/R(t+i0)-1/R(t-i0)=-2i/r(t)$ on $(a,b)$ and $+2i/r(t)$ on
$(-b,-a)$ by \cref{lem:R}. Then
\[
\frac1{R(z)}=\frac1\pi\int_a^b\frac{\dd t}{r(t)}\Big(\frac1{z-t}-\frac1{z+t}\Big).
\]
At $z=x+i0$ the left side is $-i/r(x)$, purely imaginary, so the real part
of the right side vanishes:
$\PV\int_a^b\frac{\dd t}{r(t)(x-t)}=\int_a^b\frac{\dd t}{r(t)(x+t)}$. Since
$\frac{2x}{x^2-t^2}=\frac1{x-t}+\frac1{x+t}$, the second form equals the
third.
\end{proof}

\subsection{The mean of $\alpha_\star$}

By Fubini's theorem applied to the first form,
\begin{align*}
\int_a^b\alpha_\star\dd x&=\tau_\star\int_a^b\frac{(m^2-s^2)(b-s)}{r(s)}\dd s\\
&=\tau_\star\big(b(m^2I_0-I_2)-m^2I_1+I_3\big),
\end{align*}
with $I_j=\int_a^bs^j\dd s/r(s)$. The first bracket vanishes by the choice of
$m$. With $\sigma=s^2$, $I_1=\frac\pi2$ and $I_3=\frac\pi4(a^2+b^2)$ by
\cref{app:integrals}, so
\begin{align*}
\int_a^b\alpha_\star&=\tau_\star\frac\pi4\big(a^2+b^2-2m^2\big)\\
&=\frac{\theta K(k)}{2b}\big(a^2+b^2\big)-\theta b\,E(k)
=\frac\theta2J,
\end{align*}
in agreement with \cref{thm:main}.\footnote{Here
$\tau_\star\pi/4=\theta K(k)/(2b)$ and $\tau_\star\pi m^2/2=\theta bE(k)$; the
last expression is $\frac\theta2\big((a^2+b^2)I_0-2I_2\big)$.}

\subsection{The lead extremal}

The system $H_\star(-ab/z)$ of \cref{thm:lead} has, in the representation
\eqref{eq:foster}, linear coefficient $0$, measure on $S$ given by its
attenuation, and an atom $ab\,\tau_\star$ at the origin (\cref{thm:lead},
proof). Its phase on the band is $-\abs\theta$. The sum rule
\eqref{eq:sumrule}, with the atom costing $\frac\pi4W(0)\,ab\,\tau_\star$ and
$W(0)=(b-a)^2/(ab)$, gives
\begin{align*}
\frac1{b-a}\int_a^b\alpha\dd\omega&=-\abs\theta\Lambda+\frac\pi4(b-a)\tau_\star\\
&=-\abs\theta\frac{K-E}\kappa+\abs\theta\kappa K\\
&=\abs\theta\,\frac{E(\kappa)-\kappa'^2K(\kappa)}\kappa ,
\end{align*}
where $K,E$ are taken at $\kappa$ and
$\frac\pi4(b-a)\tau_\star=\abs\theta(b-a)K(\kappa)/(a+b)=\abs\theta\kappa K(\kappa)$.
This is \cref{rem:leadmean}.

\section{Verification}\label{app:verification}

No statement of the paper rests on a computation: every result has a proof
in the text. Two kinds of checks accompany the source files.

\emph{Numerical checks.} The script \texttt{anc/checks.py} re-evaluates the
identities of the paper in high-precision arithmetic: the boundary values
of $W$ and the kernel identity \eqref{eq:kernel} at sample points; the
moments of \cref{lem:moments}; the three forms of $J$ and the Landen
relations; the bounds on $\Lambda$ and \cref{tab:lambda}; $m^2$, Legendre's
relation and $\tau_\star$; the three formulas of \cref{lem:alphaforms} and the
mean of $\alpha_\star$; the octave values of \cref{rem:extremal}; the lead
extremal; the single-zero reduction, its asymptotics and \cref{tab:zeros};
the ratios $J/P$ of \cref{cor:stable}; for \cref{sec:allzeros}, the root
$\eta_\star=\sqrt{ab}$ at a right angle, the agreement of the Cauchy integral
$v_{g_\sharp}$ with both closed forms \eqref{eq:alpharight}, the integral
\eqref{eq:lambdaright}, the value at the geometric centre and the symmetry,
every entry of \cref{tab:allzeros}, the positivity of $\as_\theta$ on the band for
three band ratios and three phases, the mean of $\as_\theta$ and the pointwise
comparison with $\alpha_\star$ at $\theta=1$, the narrow-band ratios after
\cref{cor:narrowzeros} and the nomes and the case of equality of
\cref{cor:nome}; for \cref{sec:pointwise}, the branch of $R$, both
identities of \cref{lem:atoms} at nine points for three bands, the
positivity of $k_t$, \eqref{eq:betarec} for a single resonance below and
above the band, \cref{lem:leadprofile} in the plane and on the band, the
integral of $r/x$, the means $\Lambda$ and $\Lambda_-$ of the two profiles,
the Herglotz property, the linear coefficient and the mean of the
compensated optimum of \cref{thm:affine}, and the octave values $A(m)$,
$\Lambda_-$ and $\vartheta_c$; the closed forms, the floor and the
decay bounds of \cref{sec:beyond}, with the Toeplitz and moment lemmas of the
construction, the constant $C_{a,b}$ and the kernels of
\cref{thm:sharperate}, and the taps and the rational bound of
\cref{prop:sixtap}; for \cref{ssec:severalexact}, the roots $c_j$, the sign and asymptotics of $W_B$, the moments and the Glauert identity of \cref{lem:glauertB}, the equality case and the positivity of $\alpha$ in \cref{thm:severalbands}, the superadditivity and the two-band numbers; and the numbers of \cref{sec:audit,sec:waves}. It reports
\checkscount\ checks passed and none failed.

\emph{Lean.} The file \texttt{lean/PhaseShiftCost.lean} is checked by the
kernel of Lean~4 (version 4.34.0, core library only). It contains
\leancount\ theorems and no \texttt{sorry}. It verifies the dimensional
analysis and the arithmetic of \cref{sec:audit}; the coefficients $d_n$ of
the series of $\Lambda$, their identification with
$\binom{2n}n^24^{-2n}\,2n/(2n-1)$ for $n\le60$, and their monotonicity for
every $n$; enclosures of $\Lambda/\pi$ to nine decimals for six band ratios,
from partial sums of the series and the coefficient bound behind the tail
estimate; and the polynomial identities used in
\cref{sec:test,sec:constant,sec:sumrule,sec:zeros,sec:allzeros} and in the
appendices; for \cref{sec:allzeros} these are the identities behind
\eqref{eq:Hright}, both forms of \eqref{eq:alpharight}, the symmetry
$\omega\mapsto ab/\omega$, the value at $\sqrt{ab}$ and \eqref{eq:lambdaright},
and for \cref{sec:pointwise} the signs of the kernels, the comparison of
\cref{cor:compare}, the step of \cref{cor:pointwisecons}(iv), the vanishing
linear coefficients of the compensated optima and the threshold
$\vartheta_c$, and for \cref{sec:beyond} the identities and constants of
\cref{prop:floor,thm:expdecay,cor:rate}, the determinants of
\cref{thm:noloewner}, and the polynomial identities and the rational bound
$\int_1^2\alpha<7/10000$ of \cref{prop:sixtap}, computed exactly from the
partial sums of the cosine series. Of the \leancount\ theorems, \leanfree\ depend on no axiom, two on
\texttt{propext} only, three on \texttt{propext} and \texttt{Quot.sound}, and the others on Lean's three standard axioms
(\texttt{propext}, \texttt{Classical.choice}, \texttt{Quot.sound}), through
the tactics \texttt{grind} and \texttt{omega}.

The analytic theorems are not formalised. The sum rule
(\cref{thm:sumrule}), the identification of $\Lambda$ with elliptic integrals
and with its series, the extremal system, the reconstruction of
\cref{thm:reconstruction}, and the results on zeros rest on
boundary values of Cauchy integrals, the M.~Riesz theorem, Hardy spaces,
conformal maps and extremal length, Landen's
transformation and the summation of a geometric series, none of which is
formalised in the file. They are proved in the text and in
\cref{app:elliptic,app:integrals,app:profiles}.
